\section{\texorpdfstring{Coupling independence for \(q\ge1.809\,\Deg\), \(\Deg\ge125\)}
{Coupling independence for q >= 1.809 Delta, Delta >= 125}}
\label{sec:below-vigoda}

Set
\begin{equation}\label{eq:cv-main-constants}
  m_0:=\frac{1724}{1809},
  \qquad
  \delta_{\rm CV}:=\frac{59}{226125}.
\end{equation}

\begin{theorem}[Hamming coupling independence at $1.809\,\Deg$]
\label{thm:cv-ci}
Let \(q,\Deg\) be positive integers with \(\Deg\ge125\) and
\(q\ge1.809\,\Deg\).  For every \(G\in\Gdeg\), every pinning \(\tau\),
every \(x\in[0,1]\), every \(r\in V^\tau\), and all
\(a,b\in\colours\),
\begin{equation}\label{eq:cv-ci}
  \WHam\!\left(
    \mupin{G}{\tau^{r=a}}{x},
    \mupin{G}{\tau^{r=b}}{x}
  \right)
  \le \frac{2}{m_0\delta_{\rm CV}}
  =\frac{409060125}{50858}<8043.19.
\end{equation}
The laws live on \(V^\tau\setminus\{r\}\).  At \(x=0\), they are the
normalized hard pinned laws.
\end{theorem}

At \(x=0\), the underlying list-colouring contraction is due to Carlson and
Vigoda~\cite[Theorem~4.12 and the proof of Corollary~1.2]{CarlsonVigoda2024}.
Combined with the stationary-law comparison of
\cite[Lemma~4.3]{BlancaEtAl2022Coupling}, it already yields the corresponding
bounded conditional Hamming couplings; the explicit coupling-independence
formulation appears in \cite[Theorem~20 and Proposition~21]{CFFGZZ2024}.
This section extends that input to positive activities.

The proof has three stages.  Random active constraints first turn the soft
Potts law into a mixture of hard list-colouring laws, on which we use the
Carlson--Vigoda profile.  A coupling matrix for adjacent configurations
then combines the Carlson--Vigoda component matches with a product
completion, and a geometric path metric gives uniform contraction.
Finally, stationary comparison through the root-deleted instance proves
\Cref{thm:cv-ci}; the low-multiplicity estimate is certified exactly in
\Cref{lem:cv-certificate}.

\subsection{Soft Carlson--Vigoda kernel}

The aggregate flip parameters of~\cite[Eq.~(1)]{CarlsonVigoda2024} are
\begin{equation}\label{eq:cv-params}
\begin{aligned}
  &P_0=0,\quad P_1=1,\quad P_2=\frac{81}{250},\quad
    P_3=\frac{77}{500},\\
  &P_4=\frac{11}{125},\quad P_5=\frac{11}{250},\quad
    P_6=\frac{11}{1000},\quad P_j=0\quad(j\ge7),
\end{aligned}
\end{equation}
and we write \(\boldsymbol P:=(P_j)_{j\ge0}\).  A feasible alternating
component of size \(j\) is accepted with probability \(P_j/j\) from each
of its \(j\) proposal vertices, so the aggregate move has transition
probability \(P_j/(nq)\) on an \(n\)-vertex instance.  We shall use
\[
  \max_j jP_j=1,
  \qquad
  \max_j(j-2)P_j=\frac{22}{125}.
\]

Throughout the contraction proof, fix \(G\in\Gdeg\), a pinning
\(\tau\), \(x\in(0,1]\), and \(n:=\lvert V^\tau\rvert\ge1\); assume
\(\Deg\ge125\), \(q\ge(1809/1000)\Deg\), and set \(\theta:=1-x\).
We use the free graph \(G^\tau\) and boundary counts \(\bcount_u^\tau(c)\)
from \Cref{sec:potts-setup}, which satisfy
\[
  \deg_{G^\tau}(u)+\sum_{c\in\colours}\bcount_u^\tau(c)
  =\deg_G(u)\le\Deg.
\]
The free--pinned constraints retain their individual edge labels, as in
\Cref{sec:soft-flip-interface}, so repeated boundary colours are counted
with multiplicity.

\begin{definition}[Soft-CV kernel]\label{def:cv-kernel}
For \(G\in\Gdeg\), a pinning \(\tau\), and \(x\in(0,1]\), define
\[
  K_{x,\mathrm{CV}}^{G,\tau}:=K_{x,\boldsymbol P}^{G,\tau}.
\]
This is the soft profile kernel of \Cref{sec:soft-flip-interface}
specialized to \eqref{eq:cv-params}.
\end{definition}

The endpoint $x=0$ is added in \Cref{thm:cv-ci} by convergence of the
finite-state child laws.

Because component flips are involutions and $P_j/j$ compensates the $j$
proposal vertices, the profile-parametrized argument of
\Cref{lem:soft-stationary} shows that \(K_{x,\mathrm{CV}}^{G,\tau}\) is irreducible and
reversible with stationary law $\mupin{G}{\tau}{x}$.  For adjacent
configurations we use the coupled activations of
\Cref{lem:coupled-activation-root-local}; in each marginal the subsequent
move is the ordinary hard list-colouring flip transition.

\subsection{Geometric path metric}

Let \(X,Y\in\colours^{V^\tau}\) be adjacent with disagreement root \(v\),
\(X_v=a\), \(Y_v=b\),
$T=\set{a,b}$. A neighbour $u\in N_{G^\tau}(v)$ is \emph{regular} if
$X_u=Y_u\notin T$. For regular $u$ define
\begin{equation}\label{eq:cv-scores}
  r_u:=\abs{\set{y\in N_{G^\tau}(u)\setminus\set v:X_y\in T}},\qquad
  e_u:=\bcount_u^\tau(a)+\bcount_u^\tau(b),\qquad
  s_u:=\theta\,x^{\,r_u+e_u}.
\end{equation}
The survival probability $s_u$ is exactly the probability that, in the
active instance,
the root edge $vu$ is active, every free $T$-blocker at $u$ is inactive,
and neither target colour is deleted at $u$.  Thus $s_u$ is a
deterministic function of the input pair; the probability is
over the subsequently sampled active set. Expectations below average both the
active-set draw and the flip transition.
Write \(\mathcal U_u\) for this activation event, so that
\(\Pr(\mathcal U_u)=s_u\).

\begin{definition}[Geometric soft-CV metric]\label{def:cv-metric}
For adjacent $X,Y$ as above set
\begin{equation}\label{eq:cv-edge-length}
  \ell_x^{G,\tau}(X,Y):=1-\frac{17}{200\,q}\sum_{u\ \text{regular}}s_u ,
\end{equation}
and let $d_x^{G,\tau}$ be the shortest-path metric generated by these
one-coordinate edge lengths on the graph with vertex set
\(\colours^{V^\tau}\), in which two configurations are adjacent exactly when
they differ at one vertex. The
coefficient is $(P_2-P_3)/2$ per unit of $q$, matching the hard metric
of~\cite[Eq.~(2)]{CarlsonVigoda2024}. Its limit as $x\downarrow0$ is the
hard metric: a neighbour has limiting survival probability one exactly when it has no blocker
and no target deletion, and otherwise has limiting survival probability zero.
\end{definition}

\begin{lemma}[Metric comparison; adjacent geodesics]\label{lem:cv-metric}
If $q\ge1.809\,\Deg$, then
$m_0\Ham(\sigma,\sigma')\le d_x^{G,\tau}(\sigma,\sigma')\le\Ham(\sigma,\sigma')$ for
all \(\sigma,\sigma'\in\colours^{V^\tau}\), where \(m_0\) is given in
\eqref{eq:cv-main-constants}; for adjacent pairs the direct edge is a geodesic.
\end{lemma}

\begin{proof}
There are at most $\Deg$ regular neighbours and each survival probability is at most $1$,
so every edge length lies in
$[1-17\Deg/(200q),1]\subseteq[1724/1809,1]$. Any path makes at least
$\Ham(\sigma,\sigma')$ coordinate changes, and changing disagreements one at a
time costs at most $1$ each. A path avoiding the direct adjacent edge uses
at least two edges, of cost at least $2m_0>1$.
\end{proof}

For \(\sigma,\sigma'\in\colours^{V^\tau}\) and a disagreement \(y\), put
\(T_y=\set{\sigma_y,\sigma'_y}\).  Call \(w\) an
\emph{endpoint-common neighbour} of \(y\) when
\[
 w\in N_{G^\tau}(y),\qquad \sigma_w=\sigma'_w\notin T_y.
\]
For such \(y,w\), define the union blocker count
\[
\begin{aligned}
b_{\sigma,\sigma'}(y,w)
 :={}&\abs{\set{s\in N_{G^\tau}(w)\setminus\set y:
          \sigma_s\in T_y\ \text{or}\ \sigma'_s\in T_y}}\\
 &+\sum_{c\in T_y}\bcount_w^\tau(c).
\end{aligned}
\]
Thus each free-edge factor is counted once even if it is a blocker in
both endpoint configurations, and the edge \(yw\) is excluded.  Define
\begin{equation}\label{eq:cv-output-score}
 \mathcal S_{\rm out}(\sigma,\sigma')
 :=\sum_{y:\sigma_y\ne\sigma'_y}
      \theta\!\!\sum_{\substack{w\in N_{G^\tau}(y)\\
                 \sigma_w=\sigma'_w\notin T_y}}
        x^{b_{\sigma,\sigma'}(y,w)} .
\end{equation}

\begin{lemma}[A path upper bound for the geometric metric]\label{lem:cv-path}
For all \(\sigma,\sigma'\in\colours^{V^\tau}\),
\begin{equation}\label{eq:cv-path}
  d_x^{G,\tau}(\sigma,\sigma')\le\Ham(\sigma,\sigma')
   -\frac{17}{200q}\mathcal S_{\rm out}(\sigma,\sigma').
\end{equation}
\end{lemma}

\begin{proof}
Change the disagreements in a fixed order. At the step changing $y$, every
blocker present in the intermediate configuration carries a colour taken
from $\sigma$ or from $\sigma'$, so its blocker count at $w$ is at most
$b_{\sigma,\sigma'}(y,w)$.  Since $m\mapsto x^m$ is nonincreasing for
$x\in(0,1]$, each step's edge length is at most
$1-\tfrac{17}{200q}\theta\sum_wx^{b_{\sigma,\sigma'}(y,w)}$. Sum along the
path.
\end{proof}

\subsection{Conditional coupling under fixed activations}

\subsubsection{Common off-root components}
Fix adjacent \(X,Y\in\colours^{V^\tau}\) and an outcome of the coupled activations from
\Cref{lem:coupled-activation-root-local}.  Write the two active hard
instances as
\[
  F_X=(V^\tau,E_X,L_X),\qquad F_Y=(V^\tau,E_Y,L_Y).
\]
Here \(E_X,E_Y\) are the active free-edge sets, and
\(L_X=(L_{X,u})_{u\in V^\tau}\), \(L_Y=(L_{Y,u})_{u\in V^\tau}\)
are the corresponding effective-list families.  The two instances agree
away from $v$.  Call every \(c\notin T\) a
\emph{regular colour}, and put
\[
  \Gamma_c:=\set{u\in N_{G^\tau}(v):
      X_u=Y_u=c,\ vu\in E_X\cap E_Y},\qquad k_c:=\abs{\Gamma_c}.
\]
Lemma~\ref{lem:coupled-activation-root-local} implies that an edge $vu$ with common colour
$c\notin T$ is active on one side if and only if it is active on the other.
For such a regular colour, say that \(c\) is \emph{root-allowed} when
\[
  c\in L_{X,v}\cap L_{Y,v}.
\]
The two membership tests are equivalent because the root lists can differ
only in the colours \(a,b\).  Write
\[
  \mathsf A_c:=\one{c\text{ is root-allowed}}.
\]
For a root target \(c_0\in T\), root-allowed
means \(a\in L_{Y,v}\) when \(c_0=a\), and \(b\in L_{X,v}\) when \(c_0=b\).
Equivalently it again means \(c_0\in L_{X,v}\cap L_{Y,v}\), since the other
starting configuration already uses \(c_0\).

\begin{lemma}[Components near the disagreement vertex]\label{lem:cv-root-local-structure}
For each regular colour $c\notin T$:
\begin{enumerate}
\item Deleting $v$ from the root-containing $\set{a,c}$-component in
  $F_X$ gives exactly the off-root $\set{a,c}$-components in $F_Y$ that
  meet $\Gamma_c$, with the same post-swap assignments.  The analogous
  statement holds with $a,X,Y$ replaced by $b,Y,X$.  Consequently a
  root-containing move is feasible exactly when $c$ is root-allowed
  and all the corresponding off-root moves are feasible.
\item For a fixed off-root proposal vertex and a fixed unordered colour
  pair, the components grown on the two sides have the same vertex set and
  the same post-swap feasibility whenever both avoid $v$.  If they differ,
  at least one contains $v$, and its flip recolours $v$.
\end{enumerate}
\end{lemma}

\begin{proof}
By \Cref{lem:coupled-activation-root-local}, every off-root edge and list
is the same in $F_X$ and $F_Y$.  A one-sided root edge has neighbour
colour $a$ or $b$, and a one-sided root deletion deletes $a$ or $b$.
Thus all edges and lists relevant to a colour $c\notin T$ are common away
from the root.

Hence the $\set{a,c}$-component through $v$ on the
$X$-side is $v$ together with exactly the off-root
$\set{a,c}$-components attached through $\Gamma_c$. These are the same
components on the $Y$-side, because changing the colour at $v$
from $a$ to $b$ does not change the induced $\set{a,c}$ graph away from
$v$. The post-swap assignments agree at every off-root vertex. At $v$ the
root-containing swap assigns $c$; the possible extra deletion of the old
colour $a$ is therefore harmless. Thus the only additional root list check
is that $c$ is root-allowed, proving part~1. The same off-root identity
proves part~2: the component grown from an off-root proposal vertex can first differ
only when one growth reaches $v$, and any nontrivial two-colour flip
containing $v$ changes its colour.
\end{proof}

\begin{lemma}[Partition of the aggregate moves]\label{lem:cv-move-partition}
After common off-root aggregate flips are coupled synchronously, every
remaining nonidentity aggregate move belongs
uniquely to one of the regular-colour families \(c\notin T\) or to one of
the two root-colour families \(a,b\).

More explicitly, delete \(v\) and consider the common off-root
\(\{a,b\}\)-components.  Let \(\mathcal B_a\) be those containing an
active \(a\)-coloured neighbour joined to \(v\) on the \(Y\)-side, and let
\(\mathcal B_b\) be those containing an active \(b\)-coloured neighbour
joined to \(v\) on the \(X\)-side.  Then
\[
\begin{array}{c|c}
\text{block class}&\text{assigned moves}\\ \hline
\mathcal B_a\setminus\mathcal B_b&
  \text{\(a\)-family: \(Y\)-root and \(X\)-off-root}\\
\mathcal B_b\setminus\mathcal B_a&
  \text{\(b\)-family: \(X\)-root and \(Y\)-off-root}\\
\mathcal B_a\cap\mathcal B_b&
  \text{contained in the two root moves, with no off-root copy}\\
\text{neither}&\text{synchronously coupled}.
\end{array}
\]
\end{lemma}

\begin{proof}
When the corresponding components on both sides avoid \(v\), they have
the same flipped set, feasibility, and post-swap assignments on that set
by \Cref{lem:cv-root-local-structure}; these moves are coupled synchronously.
Every remaining move either contains \(v\) or is an off-root counterpart
of a component that contains \(v\).  The root-containing component has
colour pair \(\{a,c\}\) on the \(X\)-side or \(\{b,c\}\) on the
\(Y\)-side.  If \(c\notin T\), this determines the regular \(c\)-family
uniquely.

It remains to consider the pair \(\{a,b\}\).  Away from \(v\), its
components are common to \(F_X\) and \(F_Y\).  A block reaches \(v\) on
the \(Y\)-side exactly when it belongs to \(\mathcal B_a\), and reaches
\(v\) on the \(X\)-side exactly when it belongs to \(\mathcal B_b\).
The four cases in the table follow.  In particular, a block meeting both
types of root neighbour is part of both root-containing components and is
not also counted as an off-root move.  This proves both exhaustion and
uniqueness.
\end{proof}

\subsubsection{Component matching and coupling completion}

\begin{definition}[Conditional coupling matrix]
\label{def:cv-component-matching}
Let \(\mathcal M_X\) and \(\mathcal M_Y\) be the side-labelled feasible
component moves from \(X\) and \(Y\), each augmented by a holding move
\(\varnothing\).  For \(M\in\mathcal M_X\) and
\(N\in\mathcal M_Y\), write \(X^M,Y^N\) for the outputs and
\(\abs M,\abs N\) for the numbers of flipped vertices, with
\(X^\varnothing=X\), \(Y^\varnothing=Y\), and
\(\abs\varnothing=0\).  Denote the two move probabilities by
\(\kappa_X(M)\) and \(\kappa_Y(N)\).  In particular, a feasible
nonholding move of size \(j\) has probability \(P_j/(nq)\).

A partial coupling matrix is a nonnegative array
\(\boldsymbol\pi=(\pi_{M,N})\) satisfying
\[
  \sum_N\pi_{M,N}\le\kappa_X(M),
  \qquad
  \sum_M\pi_{M,N}\le\kappa_Y(N).
\]
Initialize it at zero.  Match every common off-root move with its
counterpart, with the full probability on that entry: the flipped sets
and their post-swap assignments agree, but the full configurations still
differ at \(v\).  Leave the holding row and column unassigned.  The remaining entries are the
aggregate-parameter greedy
matches of~\cite[\S5.1]{CarlsonVigoda2024}, specified as follows.

Fix \(c\notin T\).  If \(k_c=0\) and \(c\) is root-allowed, put the full
probability \(1/(nq)\) of the two singleton root moves in their common
entry; their outputs agree at \(v\).  If \(c\) is not root-allowed,
neither move is feasible.  For \(k_c\ge1\), fix an order of
\(\Gamma_c\) and assign every distinct off-root component to the first
vertex of \(\Gamma_c\) that it contains.

If the \(X\)-side root-containing \(\{a,c\}\)-move is feasible, put its
full probability in the entry joining it to a largest opposite off-root
\(\{a,c\}\)-move.  Do the same for the \(Y\)-side root-containing
\(\{b,c\}\)-move.  For each family, let its maximum-index set consist of
the indices to which a largest component is assigned.  If the two sets
intersect, choose both entries at one common index; otherwise choose any
maximizer in each family.  Root feasibility implies feasibility of every
constituent off-root move by \Cref{lem:cv-root-local-structure}, and the
profile is nonincreasing on positive component sizes, so neither entry
overfills its opposite row or
column.

Subtract these entries, scan \(\Gamma_c\) in the fixed order, and at each
vertex put the largest possible entry between the two current residual
component probabilities assigned there.  Each subtraction is made from the
component's global residual, so a component meeting several vertices is
never reused.  The two flipped sets in such an entry share the indexing
vertex and therefore save at least one unit relative to their separate size
bounds.

For root colour \(a\), the root-containing move is the \(Y\)-side
\(\{a,b\}\)-move through \(v\), and its possible opposite moves are the
\(X\)-side off-root \(\{a,b\}\)-components meeting an active
\(a\)-coloured root neighbour but not \(v\); interchange \(X,Y\) for
root colour \(b\).  If there is exactly one such neighbour and both moves
are feasible, put the full root-move probability in the root/off-root
entry.  Otherwise put a feasible root move in the entry with holding on the
opposite side.  That entry removes the original root disagreement.
\end{definition}

\begin{lemma}[Component matching and product completion]
\label{lem:cv-completion}
The entries in \Cref{def:cv-component-matching} extend to a coupling matrix
with row sums \(\kappa_X\) and column sums \(\kappa_Y\).
\end{lemma}

\begin{proof}
The regular-colour capacity checks are part of the construction.  In
particular, a component is recorded at its first, or canonical, vertex of
\(\Gamma_c\), and every later match subtracts from its global residual.
This covers the case in which one component meets several vertices while
the components on the other side are distinct.  When the construction uses
a root/off-root entry for a root colour with one active neighbour, its
eligible partner \(C\) satisfies that the root component is
\(\{v\}\cup C\).  Root feasibility implies feasibility of the
off-root move and \(P_{\abs C+1}\le P_{\abs C}\), so the root/off-root
entry fits.  In every other root-colour case, at most one root move on each
side is joined to holding.  The proposals whose proposed colour is the
current colour give each side holding probability at least \(1/q\), while
a root move has probability at most \(1/(nq)\).  Hence the holding entries
also fit.

Let the residual row and column probabilities be
\begin{equation}\label{eq:cv-residual-masses}
  \alpha_M:=\kappa_X(M)-\sum_N\pi_{M,N},
  \qquad
  \beta_N:=\kappa_Y(N)-\sum_M\pi_{M,N}.
\end{equation}
They are nonnegative by the preceding checks.  Both move distributions
have total mass one, and each specified entry removes the same mass from a
row and a column, so
\[
  \sum_M\alpha_M
  =1-\sum_{M,N}\pi_{M,N}
  =\sum_N\beta_N
  =:R.
\]
If \(R>0\), add
\[
  \frac{\alpha_M\beta_N}{R}
\]
to every entry \((M,N)\); if \(R=0\), nothing remains.  This product
completion has residual row sums \(\alpha_M\) and column sums
\(\beta_N\), and therefore gives the required coupling matrix.
\end{proof}

\begin{lemma}[Cost of the completed coupling]
\label{lem:cv-completion-cost}
For the residual masses in \eqref{eq:cv-residual-masses}, the
product-completion entries have total Hamming increment at most
\[
  \sum_M\alpha_M\abs M+\sum_N\beta_N\abs N.
\]
An entry matching a common off-root move has zero increment.  An explicit
root-colour move paired with holding saves two units relative to the
flipped-set-size bound.
Every nonidentity move occurs in exactly one component family.
\end{lemma}

\begin{proof}
For every pair of moves,
\[
  \Ham(X^M,Y^N)-1\le\abs M+\abs N,
\]
because away from the two flipped sets only the original disagreement can
remain.  Consequently the product-completion cost is at most
\[
  \sum_M\alpha_M\abs M+\sum_N\beta_N\abs N.
\]
This proves the displayed residual bound.  A common off-root match leaves
only the original disagreement and hence has zero increment.  An explicit
root-colour move paired with holding changes the root to its colour on the
opposite side.  For such an $X$-side move $M$,
\(\Ham(X^M,Y)-1\le |M|-2\), and symmetrically on the other side.
This is the stated two-unit saving relative to the size bound.  Finally,
\Cref{lem:cv-move-partition} ensures that every nonidentity move occurs in
exactly one family.  Thus, even when a product-completion entry joins moves
from different families, its two size terms are charged separately to the
unique families of its row move and column move; no residual mass is omitted
or counted twice.
\end{proof}

For the rest of this section, let \((X_1,Y_1)\) denote the outputs drawn
from the completed coupling matrix.

For every \(c\notin T\), define its physical root multiplicity by
\[
 m_c:=\abs{\set{u\in N_{G^\tau}(v):X_u=Y_u=c}},
\]
and recall that \(k_c=\abs{\Gamma_c}\) is its active root multiplicity.
For the current active outcome, put
\[
 \mathcal C_{\rm low}:=
    \set{c\notin T:c\text{ is root-allowed},\
                      m_c\in\set{1,2}},
\]
and let \(E_v\) be the event that at least one selected component move
recolours \(v\).

\begin{lemma}[Conditional component budgets]\label{lem:cv-component-budgets}
Fix a coupled active outcome for adjacent $X,Y$, with disagreement root
$v$, $X_v=a$, $Y_v=b$, and put $T=\set{a,b}$. The following facts hold
for the component coupling above.
\begin{enumerate}
\item \emph{Root-event budget.}
  Conditionally on the active outcome,
  \begin{equation}\label{eq:cv-root-recolour-budget}
    nq\,\Pr(E_v\mid F_X,F_Y)
       \le q-P_2\sum_{c\in\mathcal C_{\rm low}}k_c .
  \end{equation}
\item \emph{Output-discount bound.}
  Fix a regular neighbour $u$ and
  $w\in N_{G^\tau}(u)\setminus\set v$ which is not initially
  $T$-coloured. For each $t\in T$, collect the active free edges from
  $w$ to a $t$-coloured vertex and the active pinned-neighbour constraints at $w$
  deleting $t$. After root-recolouring moves have been charged to
  $E_v$, the total transition probability that can recolour $w$ into a colour
  of $T$ is at most $(1+P_2)/(nq)$ whenever at least one collected constraint is
  active. If none is active and the root edge $vu$ is active while all
  off-root $T$-blockers and target deletions at $u$ are inactive, the
  same-neighbour matrix entry is at least $(1-P_2)/(nq)$ and flips only $u$ on
  either side. Credits obtained this way for different ordered pairs
  $(u,w)$ refer to different summands of \eqref{eq:cv-path}.
\end{enumerate}
\end{lemma}

\begin{proof}
For part~1, assign the \(X\)-side \(\{a,c\}\) and \(Y\)-side
\(\{b,c\}\) root-containing flips to the regular colour \(c\notin T\);
the two remaining \(\{a,b\}\) root flips form the \(a\)- and
\(b\)-families.  These assignments are disjoint.  A
colour that is not root-allowed contributes no root-changing move, except possibly
$a$ or $b$; for those exceptional colours the extra proposal is the
current colour on its witness side and hence is a stay proposal, while it
is forbidden on the other side. Thus it contributes zero. For
$c\notin T$ with $k_c=0$, the two root-containing components are both the
singleton $\set v$ and are coupled together with probability $1/(nq)$.
If $k_c\ge1$, each root-containing component has size at least
$k_c+1$, so their combined probability is at most $2P_{k_c+1}/(nq)$. When
$k_c=1$,
\[
  2P_2\le1-P_2
  \qquad\text{because}\qquad 3P_2=\frac{243}{250}\le1,
\]
and when $k_c=2$,
\[
  2P_3\le1-2P_2
  \qquad\text{because}\qquad
  2P_3+2P_2=\frac{239}{250}\le1 .
\]
Hence every $c\in\mathcal C_{\rm low}$ saves at least $P_2k_c$ from its
unit baseline. All other regular colours cost at most their unit
baseline. Assign the $X$-side root-containing $\set{a,b}$-move to the
$b$-baseline and the $Y$-side root-containing move to the $a$-baseline.
Each has probability at most
$P_1/(nq)=1/(nq)$, so summing the $q$ baselines and multiplying by $nq$ proves
\eqref{eq:cv-root-recolour-budget}.

For part~2, fix a target $t\in T$ and compare the two
$\set{X_w,t}$-components from $w$. If both avoid $v$,
by \Cref{lem:cv-root-local-structure} they are identical and hence
synchronously coupled, so they contribute only one copy of their component
probability. If a component contains $v$, its move is already charged to
$E_v$; after removing all such moves, again at most one off-root component
remains. Its probability is at most $\max_jP_j/(nq)=1/(nq)$. If an active free edge
joins $w$ to an off-root $t$-coloured vertex, that component has size at
least two and its probability is at most $P_2/(nq)$; if an active pinned-neighbour constraint deletes
$t$ at $w$, the component is infeasible. In either case one target
contributes at most $P_2/(nq)$ and the other at most $1/(nq)$.
The remaining possibility is that the active constraint is the edge
$wv$. Then $w$ is itself a regular root neighbour. After the two
root-containing moves are assigned to $E_v$, the possible off-root moves
are the two opposite-target components through $w$. If both are
singletons, let $r_1,r_2$ be their remaining probabilities multiplied by
$nq$, after the root/off-root entries have been assigned, taking zero for
an infeasible move. Since each original move has probability at most
$1/(nq)$, we have $0\le r_1,r_2\le1$.
The same-neighbour rule couples residual mass $\min\{r_1,r_2\}/(nq)$.
Hence, outside $E_v$, the probability that at least one of these moves
occurs is at most
\[
 \frac{r_1+r_2-\min\{r_1,r_2\}}{nq}
 =\frac{\max\{r_1,r_2\}}{nq}\le\frac1{nq}.
\]
If either is not a singleton, its probability is
at most $P_2/(nq)$, while the other is at most $1/(nq)$. Thus this case also
has total probability at most $(1+P_2)/(nq)$, proving the assertion for every active
collected constraint.

Under the safe event in the last assertion, the two target components at
$u$ are feasible singletons. Each has probability $P_1/(nq)=1/(nq)$. A
root-containing component meeting the active edge $vu$ has size at least two and consumes
at most $P_2/(nq)$ of the corresponding singleton in the root/off-root
entry. The same-neighbour entry therefore has probability at least $(1-P_2)/(nq)$ on both
sides. It flips only $u$, to opposite colours in $T$, and leaves every
other vertex unchanged. Thus its output correction in
\eqref{eq:cv-path} is indexed by the ordered pair $(u,w)$. Different
ordered pairs are different terms of the double sum there (even if they
share one endpoint), so summing these credits cannot reuse one metric
summand. The same coupled move may realize several such savings
simultaneously, in which case its path-metric correction is exactly the sum
of the corresponding $(u,w)$-summands.
\end{proof}

\subsection{Geometric discount estimates}

Recolouring a unique blocker to a fresh colour creates discount; recolouring
the root, the regular neighbour, or an off-root neighbour can destroy it.  The following
three lemmas bound these effects.

Put
\begin{equation}\label{eq:cv-alpha}
  \alpha_{\mathrm{fr}}:=q-\Deg-2 .
\end{equation}
For a regular neighbour \(u\), let \(\mathsf{SB}_u\) be the activation
event that \(vu\) is active, exactly one of the \(r_u\) other free
\(T\)-blocker edges at \(u\) is active, and every remaining \(T\)-blocker
edge and \(T\)-coloured pinned-neighbour constraint at \(u\) is inactive.
Write
\[
  b_u^{\rm SB}:=\Pr(\mathsf{SB}_u),
\]
where the probability is over the active-constraint draw.

\begin{lemma}[Fresh-colour gain]\label{lem:cv-fresh}
Fix a regular neighbour $u$ with $R=r_u$ free target blockers and \(e=e_u\) target
pinned-neighbour constraints. The expected increase of the output $u$-discount relative
to its deterministic input value $s_u$, caused by singleton recolourings of
its blockers, is at least
\[
  \frac{\alpha_{\mathrm{fr}}}{nq}\,
  b_u^{\rm SB},
\]
and
\[
  b_u^{\rm SB}=
  \begin{cases}
    R\,\theta^2x^{R+e-1},&R\ge1,\\
    0,&R=0.
  \end{cases}
\]
\end{lemma}

\begin{proof}
The displayed probability is over the active-set draw: the root edge is
active, exactly one blocker edge is active, and all other blockers and all
\(e\) target-deleting pinned-neighbour constraints are inactive --- independent activation bits
of distinct constraints.  Equivalently, the factor
\(R\theta^2x^{R+e-1}\) is the sum, over the \(R\) possible unique blocker
incidences, of the deterministic discount increment available on the
corresponding activation event; it is not an additional probability factor
for the singleton flip below.  Fix a blocker incidence $(u,s)$. At $s$, exclude the
two target colours, every colour of
its full free neighbourhood, and every colour forbidden by one of its
pinned-neighbour constraints: at least
\[
  q-2-\left(\deg_{G^\tau}(s)+\sum_{c\in\colours}\bcount_s^\tau(c)\right)
  \ge\alpha_{\mathrm{fr}}
\]
colours remain, and for each the alternating component of $s$ is a list-feasible
singleton in \emph{every} active outcome, accepted with probability $P_1=1$
and transition mass $1/(nq)$. The flip removes one blocker incidence and raises the
corresponding output $u$-discount relative to $s_u$ by
$\theta x^{R+e-1}-\theta x^{R+e}=\theta^2x^{R+e-1}$. Summing over
the $\alpha_{\mathrm{fr}}$ colours and $R$ incidences gives the claim; if one blocker
vertex serves several regular neighbours, the single flip increases the total
discount by the sum of the individual increments. Incidence summation
therefore counts each discount contribution, while the flip probability
is counted once.
\end{proof}

For a regular root-neighbour colour \(c\), recall that \(m_c\) is its
physical multiplicity, and recall that \(\bcount_v^\tau(c)\) counts the root
pinned-neighbour constraints forbidding \(c\). Define
\begin{equation}\label{eq:cv-Lx}
  L_x:=\sum_{\substack{c\notin T:\\m_c\in\set{1,2}}}
      m_c\,x^{\bcount_v^\tau(c)},\qquad
  \phi_x:=\frac{L_x}{\Deg} .
\end{equation}
Since $L_x$ is a weighted subcount of the physical root neighbours,
$0\le \phi_x\le1$.
The factor $x^{\bcount_v^\tau(c)}$ is the probability that $c$ survives in the
active root list, and root-edge activation is independent of root boundary
activation, so the expected value of
$\sum_{c\in\mathcal C_{\rm low}}k_c$ is exactly $\theta L_x$.
Averaging \eqref{eq:cv-root-recolour-budget} from
\Cref{lem:cv-component-budgets} therefore gives
\begin{equation}\label{eq:cv-root-event}
  nq\,\Pr(E_v)\;\le\;q-P_2\,\theta L_x ,
\end{equation}
where $E_v$ is the root-recolouring event. This is the soft analogue
of~\cite[Eq.~(18)]{CarlsonVigoda2024}, proved directly for the two active
hard instances.

\begin{lemma}[Ordered-incidence discount bound]\label{lem:cv-incidence}
Fix a regular neighbour $u$ and $w\in N_{G^\tau}(u)\setminus\set v$ not initially
a target blocker, so \(X_w=Y_w\notin T\).  Let \(\mathcal O_w\) be the
following disjoint set of physical constraint factors:
\[
 \mathcal O_w:=
 \set{wy\in E(G^\tau):y\ne u,\
       X_y\in T\ \text{or}\ Y_y\in T}
 \ \sqcup\
 \set{(w,t,j):t\in T,\ 1\le j\le\bcount_w^\tau(t)} ,
\]
and put \(t_w=\abs{\mathcal O_w}\). The contribution of the
ordered incidence $(u,w)$ to input discount loss minus certified output discount
is at most $(1+P_2)\,s_u/(nq)$.
\end{lemma}

\begin{proof}
Let \(\mathcal E_w\) be the event that all of \(\mathcal O_w\) is inactive. Given the
active outcome, for each target $c\in T$ the operations recolouring $w$
into $c$ flip the unique active $\set{X_w,c}$-component of $w$ in one
marginal, with transition probability at most $\max_jP_j/(nq)=1/(nq)$,
so their total probability is at most $2/(nq)$
always. When \(\mathcal E_w\) fails, some constraint of \(\mathcal O_w\) is active at
$w$: an active edge to a $c$-coloured vertex forces the
$\set{X_w,c}$-component to have size at least two, capping its transition
probability at $P_2/(nq)$, and an active pinned-neighbour constraint
forbidding $c$ makes the swap list-infeasible. Part~2 of
\Cref{lem:cv-component-budgets} shows directly that root-containing moves
are already charged to $E_v$ and common off-root components are coupled synchronously,
so the remaining probability is at most $(1+P_2)/(nq)$. This proves in the
two-instance setting the case-(i) estimate corresponding
to~\cite[\S4.3]{CarlsonVigoda2024}. The averaged probability is thus at most
\(\bigl((1+P_2)+(1-P_2)x^{t_w}\bigr)/(nq)\), and each new blocker incidence costs at most
$s_u-s_ux=\theta s_u$ of discount.

The activation bits of \(\mathcal U_u\) are disjoint from those of
\(\mathcal E_w\), so
\(\Pr(\mathcal U_u\cap\mathcal E_w)=s_u x^{t_w}\). On this event both target
components at $u$ --- the $\set{X_u,a}$-move on one side and the
$\set{X_u,b}$-move on the other --- are feasible singletons; root-containing
components meeting the active root edge have transition probability at most
$P_2/(nq)$, so after the root/off-root entry each side retains probability
at least $(1-P_2)/(nq)$, and the same-neighbour entry pairs the two moves: the coupled flip
recolours $u$ to $a$ on one side and to $b$ on the other, creating a new
disagreement at $u$ with pair $\set{a,b}$. Its blocker incidences at $w$
are exactly the \(t_w\) constraints of $\mathcal O_w$ (no other vertex
changes), so the $(u,w)$ output discount in \eqref{eq:cv-path} is
\(\theta x^{t_w}\). The certified credit is at least
\((1-P_2)\,s_u\theta x^{2t_w}/(nq)\). The net cost
is at most
\[
  \frac{s_u\,\theta}{nq}\Bigl[(1+P_2)
    +(1-P_2)\bigl(x^{t_w}-x^{2t_w}\bigr)\Bigr]
  \le\frac{(1+P_2)\,s_u}{nq},
\]
since for \(t_w\ge1\), \(x^{t_w}-x^{2t_w}\le x\) and
$(1-x)[(1+P_2)+(1-P_2)x]\le1+P_2$, with slack $x[2P_2+(1-P_2)x]\ge0$; for
\(t_w=0\) the extra term vanishes. Distinct ordered incidences certify distinct
summands of \eqref{eq:cv-path}, so no credit is reused.
\end{proof}

\begin{lemma}[Expected loss of the geometric discount]\label{lem:cv-expected}
Fix the adjacent input pair and a deterministic order of the off-root
neighbours of each regular \(u\).  For each coupled outcome, process in order
the root recolouring, a recolouring of \(u\) into \(T\), and then the
off-root neighbours in the fixed order.  Let \(\mathsf{Loss}_u\) be the
telescoping decrease of the input \(u\)-discount caused by these operations,
charging each decrease at its first cause.  Let \(\mathsf{Credit}_u\) be the
sum of the distinct \((u,w)\)-indexed output-score summands certified by the
same-neighbour matches in \Cref{lem:cv-component-budgets,lem:cv-incidence}.
Thus every input loss is charged once and every output summand is credited at
most once.  Expectations are over both the coupled activation draw and the
coupled hard update.  For every regular neighbour \(u\),
\begin{gather*}
  \E[\mathsf{Loss}_u]-\E[\mathsf{Credit}_u]
  \;\le\;\frac{\beta^{+}_x}{nq}\,s_u ,\\
  \beta^{+}_x:=q-P_2\theta L_x+(1+P_2)\Deg+2 .
\end{gather*}
\end{lemma}

\begin{proof}
The discount at $u$ can disappear in three ways: the root is recoloured ---
probability bounded by \eqref{eq:cv-root-event}; the root survives but $u$
is recoloured into a target colour --- by
\Cref{lem:cv-root-local-structure}, common off-root components are coupled synchronously,
whereas every non-common component meeting the root is already charged to
$E_v$; after those moves are removed, each target contributes probability
at most $1/(nq)$, hence at most $2/(nq)$ in total; or a neighbour $w$ creates a new
blocker --- at most $\Deg$ ordered incidences, each bounded by
\Cref{lem:cv-incidence}. The three types can overlap on a single
operation, but each charge separately upper-bounds the part of the discount
loss it names and the total loss telescopes over the triggered types, so
summing the three expected charges is conservative.
\end{proof}

\begin{lemma}[Disjointness of the output credits]\label{lem:cv-credit-disjoint}
For every coupled outcome, the sum of all output-score summands retained as
credits in \Cref{lem:cv-fresh,lem:cv-incidence,lem:cv-expected} is at most
\(\mathcal S_{\rm out}(X_1,Y_1)\).
\end{lemma}

\begin{proof}
The ordered-incidence credit associated with \((u,w)\) is the summand of
\eqref{eq:cv-output-score} indexed by the output disagreement at \(u\) and
its endpoint-common neighbour \(w\).  The deterministic ordering in
\Cref{lem:cv-expected} retains each such ordered pair at most once.
The fresh-colour credit at a regular root neighbour \(u\) is instead the
summand indexed by the surviving root disagreement \((v,u)\).  These
summands have first coordinate \(v\), whereas every ordered-incidence
summand has first coordinate \(u\ne v\); the two collections are therefore
disjoint.  If a single coupled flip realizes several credits, they are
distinct summands in the defining double sum for
\(\mathcal S_{\rm out}\).  Their total is consequently bounded by the full
output score, outcome by outcome.
\end{proof}

\subsection{Per-colour charge estimates}

Set $\varrho=q/\Deg$, $\eta(\varrho)=\frac{P_2-P_3}{2\varrho}=\frac{17}{200\varrho}$, and
\begin{equation}\label{eq:cv-gammas}
  \gamma_{\rm gain}:=\eta(\varrho)
     \Bigl(\varrho-1-\frac2\Deg\Bigr),\qquad
  \gamma_{\rm loss}:=\eta(\varrho)
     \Bigl(\varrho-P_2\theta\phi_x+1+P_2+\frac2\Deg\Bigr),
\end{equation}
the per-$\Deg$ metric corrections supplied by
\Cref{lem:cv-fresh,lem:cv-expected}. On $\varrho\ge1809/1000$, $\Deg\ge125$,
these coefficients lie in the half-open range
\[
  (\gamma_{\rm gain},\gamma_{\rm loss})\in
  \left[\frac{13481}{361800},\frac{17}{200}\right)
  \times\left[0,\frac{799}{5400}\right],
  \qquad \frac{799}{5400}<P_2.
\]
Indeed,
\[
 \gamma_{\rm gain}=\frac{17}{200}
   \left(1-\frac{1+2/\Deg}{\varrho}\right),\qquad
 \gamma_{\rm loss}=\frac{17}{200}
   \left(1+\frac{1+P_2+2/\Deg-P_2\theta\phi_x}{\varrho}\right).
\]
Over the displayed continuous parameter domain, the first expression has
its minimum at
$(\varrho,\Deg)=(1809/1000,125)$ and approaches $17/200$ without reaching
it; the second has its maximum at
$(\varrho,\Deg,\theta\phi_x)=(1809/1000,125,0)$.  For the finite
maximization, close the open endpoint and use the
closed bounding rectangle
\begin{equation}\label{eq:cv-gamma-box}
  \mathcal R_\gamma:=
  \left[\frac{13481}{361800},\frac{17}{200}\right]
  \times\left[0,\frac{799}{5400}\right].
\end{equation}
For any coefficient pair in this rectangle, set
\begin{equation}\label{eq:cv-lambda-rates}
\begin{aligned}
 \Lambda_1&:=2-P_2+\gamma_{\rm loss},&
 \Lambda_2&:=2-P_3-\gamma_{\rm gain},\\
 \Lambda_3&:=2-P_2+2P_3-2P_4,&
 \Lambda_{\rm low}&:=\max\set{\Lambda_1,\Lambda_2,\Lambda_3}.
\end{aligned}
\end{equation}

\subsubsection{Finite component encoding}

\begin{definition}[Finite component records]\label{def:cv-component-records}
Fix a regular colour $c$ with $k_c\ge1$, order
$\Gamma_c=\set{u_1,\ldots,u_{k_c}}$.  At $u_i$, record the size
of the $Y$-side off-root $\set{a,c}$-component and the size of the
$X$-side off-root $\set{b,c}$-component.  A component meeting several
$u_i$ is recorded at its first listed incidence; its later appearances receive
the marker $0$, which represents no move and has $P_0=0$.

Replace every positive recorded size at least $7$ by $7$, and denote the
resulting record sizes by $a_i,b_i$.  These are truncated sizes, not
necessarily the physical sizes.  The truncation changes no transition mass:
such an off-root component has mass zero, and a root-containing component
containing it has size at least $8$ and also has mass zero.  Put
\[
  R_X=1+\sum_i a_i,\qquad R_Y=1+\sum_i b_i.
\]
 Let $\chi_i,\upsilon_i$ be the available flags of the \(a\)- and
\(b\)-records, respectively.  Let
\[
 \xi_i:=\one{\text{the \(Y\)-side \(\{a,c\}\)-flip recorded by \(a_i\)
                         is feasible}},\qquad
 \zeta_i:=\one{\text{the \(X\)-side \(\{b,c\}\)-flip recorded by \(b_i\)
                         is feasible}} .
\]
For a positive record, the availability flags are explicitly
\[
  \chi_i=\one{a\in L_{X,u_i}}=\one{a\in L_{Y,u_i}},
  \qquad
  \upsilon_i=\one{b\in L_{X,u_i}}=\one{b\in L_{Y,u_i}},
\]
where the equalities use the common off-root lists.  If \(a_i=0\), set
\(\chi_i=\xi_i=0\); if \(b_i=0\), set
\(\upsilon_i=\zeta_i=0\).  Moreover,
\[
  \xi_i\le\chi_i,\qquad \zeta_i\le\upsilon_i,\qquad
  a_i=1\Longrightarrow\xi_i=\chi_i,\qquad
  b_i=1\Longrightarrow\zeta_i=\upsilon_i.
\]
Indeed, a feasible swap must place its target colour legally at the
recording neighbour, while for a singleton this is its only post-swap list
check.  Thus every realized record satisfies these restrictions.

Let
$w_X=\mathsf A_c\prod_{i:a_i>0}\xi_i$ and
$w_Y=\mathsf A_c\prod_{i:b_i>0}\zeta_i$ be the root-component
feasibility weights.  Following \Cref{def:cv-component-matching}, choose
\[
 i_X\in\arg\max_{i:a_i>0}a_i,
 \qquad
 i_Y\in\arg\max_{i:b_i>0}b_i,
\]
using the same index whenever a common maximizer exists.  Among the
permitted pairs, choose one that maximizes the right-hand side of
\eqref{eq:cv-charge} below.  All move masses in this record are multiplied
by \(nq\); thus a feasible size-\(j\) move has scaled mass \(P_j\).
Define the \emph{record upper charge} for this arbitrary multiplicity:
\begin{equation}\label{eq:cv-charge}
  \mathcal H_c=
    w_X(R_X-a_{i_X}-1)P_{R_X}
    +w_Y(R_Y-b_{i_Y}-1)P_{R_Y}
   +\sum_i\bigl(a_i\alpha_i+b_i\beta_i-\min\set{\alpha_i,\beta_i}\bigr),
\end{equation}
where
\[
 \alpha_i=\xi_iP_{a_i}-w_X\one{i=i_X}P_{R_X},\qquad
 \beta_i=\zeta_iP_{b_i}-w_Y\one{i=i_Y}P_{R_Y}
\]
are the remaining scaled off-root masses.  The first
two terms are the root/off-root matches; after
their masses are subtracted, the sum is the residual match through the
same $u_i$.  Thus \eqref{eq:cv-charge} is the conservative record used to
bound the coupling above.  It is defined for every
$k_c\ge1$ and both values of $\mathsf A_c$.  The cases
$\mathsf A_c=1$, $k_c\in\set{1,2}$ are handled by
\Cref{lem:cv-certificate}.  Finally, define
\[
  N_{11}(c):=\abs{\set{i:a_i=b_i=1,\ \chi_i=\upsilon_i=1}},
  \qquad
  N_{12}(c):=\abs{\set{i:\set{a_i,b_i}=\set{1,2},\
                              \chi_i=\upsilon_i=1}} .
\]
The first index set uses the ordered pair $(1,1)$, whereas the second uses
$(1,2)$ or $(2,1)$.  They are therefore disjoint, and hence
\begin{equation}\label{eq:cv-small-count-budget}
  N_{11}(c)+N_{12}(c)\le k_c.
\end{equation}
A repeated appearance is never counted by $N_{11}(c)$.  Omitting a
possible $N_{12}(c)$ credit at such an appearance is conservative because
that term enters with a minus sign.
\end{definition}

For a regular colour \(c\) with \(k_c=0\), set
\[
  \mathcal H_c:=-\mathsf A_c,
  \qquad N_{11}(c):=N_{12}(c):=0.
\]

\begin{lemma}[Graph-to-record domination]\label{lem:cv-record-domination}
For every regular colour \(c\), its contribution to the conditional
expected Hamming increment, after multiplication by \(nq\), is at most
\(\mathcal H_c\).  When \(k_c=0\), this means that the contribution is at
most \(-\mathsf A_c\), which is the convention used below.
\end{lemma}

\begin{proof}
Suppose first that \(k_c\ge1\).  A root-containing component consists of
the root and its disjoint off-root components.  Matching it to the
opposite component indexed by \(i_X\), or by \(i_Y\), therefore leaves at
most \(R_X-a_{i_X}-1\), or \(R_Y-b_{i_Y}-1\), new disagreements per unit
of matched mass.  These give the first two terms of
\eqref{eq:cv-charge}.

After those masses are subtracted, the two residual moves recorded at
\(u_i\) have masses \(\alpha_i,\beta_i\) and flipped-set sizes
\(a_i,b_i\).  Their maximal match has mass
\(\min\{\alpha_i,\beta_i\}\).  Both flipped sets contain \(u_i\), so its
Hamming increment is at most
\[
 a_i\alpha_i+b_i\beta_i-\min\{\alpha_i,\beta_i\}.
\]
If the components share additional vertices, the actual increment is
smaller; this is why \(\mathcal H_c\) is an upper charge rather than an
exact cost.  Every remaining one-sided probability is passed to the
product completion and bounded by its flipped-set size.  Summing over the
canonical indices gives \eqref{eq:cv-charge}; a repeated appearance has
the zero marker and contributes no second copy.

If \(k_c=0\) and \(c\) is root-allowed, the two singleton root
moves are coupled and remove the sole input disagreement, giving
increment \(-1\).  If \(c\) is unavailable, there is no move in this
family and its contribution is \(0\).  These are exactly the two values
of \(-\mathsf A_c\).
\end{proof}

The geometric coefficients range over the rectangle
\eqref{eq:cv-gamma-box}, but the finite verification need not inspect a
subdivision of that rectangle.  The following elementary observation
reduces it to one coefficient pair.

\begin{lemma}[Balancing-point domination]\label{lem:cv-balancing-domination}
For real numbers $g_{\rm gain},g_{\rm loss}$, put
\[
 \mathcal L(g_{\rm gain},g_{\rm loss}):=
 \max\set{2-P_2+g_{\rm loss},\ 2-P_3-g_{\rm gain},\ \Lambda_3}
\]
and, for a fixed component record of a regular colour \(c\), define its
excess by
\[
 \mathcal E_c(g_{\rm gain},g_{\rm loss}):=
 \mathcal H_c+g_{\rm loss}N_{11}(c)-g_{\rm gain}N_{12}(c)
       +1-k_c\mathcal L(g_{\rm gain},g_{\rm loss}).
\]
Then
\[
 \mathcal E_c(g_{\rm gain},g_{\rm loss})
 \le \mathcal E_c\!\left(\frac{19}{500},\frac{33}{250}\right).
\]
At this balancing pair, all three terms in the maximum defining
$\mathcal L$ equal $226/125$; the pair belongs to
$\mathcal R_\gamma$ in \eqref{eq:cv-gamma-box}.
\end{lemma}

\begin{proof}
Set
\[
 u_{\rm loss}:=g_{\rm loss}-\frac{33}{250},\qquad
 u_{\rm gain}:=\frac{19}{500}-g_{\rm gain},\qquad
 M:=\max\set{0,u_{\rm loss},u_{\rm gain}}.
\]
The identities
\[
 2-P_2+\frac{33}{250}
 =2-P_3-\frac{19}{500}
 =\Lambda_3=\frac{226}{125}
\]
give
\[
 \mathcal L(g_{\rm gain},g_{\rm loss})=\frac{226}{125}+M.
\]
Consequently,
\begin{align*}
 \mathcal E_c(g_{\rm gain},g_{\rm loss})
 -\mathcal E_c\!\left(\frac{19}{500},\frac{33}{250}\right)
 &=u_{\rm loss}N_{11}(c)+u_{\rm gain}N_{12}(c)-k_cM\\
 &\le M\bigl(N_{11}(c)+N_{12}(c)-k_c\bigr)\le0,
\end{align*}
where the last inequality is \eqref{eq:cv-small-count-budget}.
\end{proof}

Two observations connect the finite expression to the geometric estimates.

\smallskip\noindent\emph{Loss charge.}
The geometric charge $\gamma_{\rm loss}$
is levied only at the $N_{11}(c)$ locations. This covers the discount
budget of \Cref{lem:cv-expected}: a location with a record of size at least two
carries an \emph{active} \(T\)-blocker, hence lies outside
\(\mathcal U_u\), the safe event defining \(s_u\), and no discount is
at risk there; and across the
availability split,
$\theta\bigl[x^{\bcount_v^\tau(c)}x^{r_u+e_u}+(1-x^{\bcount_v^\tau(c)})\bigr]\ge
\theta x^{r_u+e_u}=s_u$.  The unavailable contribution appears
separately as $(1-\mathsf A_c)k_c$ below.

\smallskip\noindent\emph{Gain credit.}
The $N_{12}(c)$ credit $\gamma_{\rm gain}$ is claimed
only at locations with exactly one active blocker edge and both targets
surviving --- an event contained in the fresh-gain event of
\Cref{lem:cv-fresh}; distinct colours claim disjoint incidences, so
the credited mass never exceeds the gain supplied.

For \(c_0\in T=\{a,b\}\), let \(C_{c_0}(F_X,F_Y)\) denote \(nq\) times
the charged Hamming increment, including the coalescence saving,
to the \(\{a,b\}\)-component family assigned to target \(c_0\) by
\Cref{def:cv-component-matching,lem:cv-completion-cost}.  Put
\[
  \mathcal H_T(F_X,F_Y):=C_a(F_X,F_Y)+C_b(F_X,F_Y).
\]

\begin{lemma}[Drift decomposition by colour]\label{lem:cv-drift-decomposition}
For coupled active instances $F_X,F_Y$ and each \(c\notin T\), retain
$\mathsf A_c=\one{c\text{ is root-allowed}}$.  Let
$\mathcal H_c(F_X,F_Y)$ be the record upper charge
\eqref{eq:cv-charge} when $k_c\ge1$, with root-containing masses set to
zero when $\mathsf A_c=0$; for $k_c=0$ put
$\mathcal H_c=-\mathsf A_c$.  With \(\mathcal H_T(F_X,F_Y)\) as defined
above, put
\[
  \mathcal D:=\E[d_x^{G,\tau}(X_1,Y_1)]-d_x^{G,\tau}(X,Y).
\]
All expectations in this statement are over the coupled activation draw
and the coupled hard transition from the fixed adjacent input pair.
Then the one-step metric drift satisfies
\begin{equation}\label{eq:cv-drift-by-colour}
\begin{aligned}
 nq\,\mathcal D\le
 \E\Bigg[\,\mathcal H_T(F_X,F_Y)
  +\sum_{c\notin T}
    \Bigl\{\mathcal H_c+
      \gamma_{\rm loss}\bigl(\mathsf A_c N_{11}(c)+
                         (1-\mathsf A_c)k_c\bigr)
      -\gamma_{\rm gain}\mathsf A_c N_{12}(c)\Bigr\}\,\Bigg].
\end{aligned}
\end{equation}
Here $N_{11}(c)=N_{12}(c)=0$ when $k_c=0$.
\end{lemma}

\begin{proof}
Put $\omega=17/(200q)$.
\Cref{lem:cv-move-partition,lem:cv-completion,lem:cv-completion-cost}
partition, complete, and bound the cost of all
aggregate moves, while \Cref{lem:cv-record-domination} bounds the actual
regular-colour Hamming contribution by its record upper charge.
Write
\[
  S_{\rm in}:=\sum_{u\ {\rm regular}}s_u,
  \qquad
  S_{\rm out}:=\mathcal S_{\rm out}(X_1,Y_1).
\]
Thus
\Cref{lem:cv-path} first gives
\begin{equation}\label{eq:cv-drift-decomp}
 nq\,\mathcal D
 \le \E\!\left[\mathcal H_T(F_X,F_Y)+
                  \sum_{c\notin T}\mathcal H_c\right]
   +nq\,\omega\bigl(S_{\rm in}-\E[S_{\rm out}]\bigr).
\end{equation}
Fix the output-discount allocation as
follows. The root, target-recolouring, and
ordered-incidence terms used in \Cref{lem:cv-expected} are allocated
first; they do not use the fresh increments of
\Cref{lem:cv-fresh}.  Allocate those fresh increments next, and discard any
remaining output discount.  This is legitimate outcome by outcome by
\Cref{lem:cv-credit-disjoint}. With
\(b_u^{\rm SB}=\Pr(\mathsf{SB}_u)\) as defined above, the two lemmas then give
\[
 S_{\rm in}-\E[S_{\rm out}]
 \le \frac{\beta_x^+}{nq}\sum_{u\ {\rm regular}}s_u
    -\frac{\alpha_{\mathrm{fr}}}{nq}\sum_{u\ {\rm regular}}b_u^{\rm SB} .
\]
By \eqref{eq:cv-gammas},
\[
 \omega\beta_x^+=\gamma_{\rm loss},\qquad
 \omega\alpha_{\mathrm{fr}}=\gamma_{\rm gain}.
\]

It remains to turn these two averaged discount bounds into the
pointwise colour terms in \eqref{eq:cv-drift-by-colour}. For a physical regular neighbour
$u$ of colour $c$, let $\iota_u$ indicate that $vu$ is active, and let
\(N^j_{c,u}\) be one exactly when \(u\) is the first recorded incidence
of its component and is counted by \(N_{1j}(c)\), and zero otherwise,
for \(j=1,2\).  Thus \(N^1_{c,u}\) contributes to \(N_{11}(c)\) and
\(N^2_{c,u}\) contributes to \(N_{12}(c)\). Root-list
survival is independent of all root-edge and blocker bits, and
\[
 \E\!\left[\mathsf A_c N^1_{c,u}
            +(1-\mathsf A_c)\iota_u\right]
 =x^{\bcount_v^\tau(c)}s_u+
   \bigl(1-x^{\bcount_v^\tau(c)}\bigr)\theta
 \ge s_u .
\]
Moreover, $\mathsf A_c N^2_{c,u}=1$ entails exactly one active target
blocker at $u$ and survival of both targets, so
  $\E[\mathsf A_c N^2_{c,u}]\le b_u^{\rm SB}$. Summing over neighbours, using
\(k_c=\sum_{u:X_u=Y_u=c}\iota_u\), and applying the two coefficient identities proves
\eqref{eq:cv-drift-by-colour}.

There is no duplicated charge or credit:
\Cref{lem:cv-move-partition} partitions the hard moves;
$\mathsf A_c$ and $1-\mathsf A_c$ are disjoint;
$N_{11}(c)$ and $N_{12}(c)$ are disjoint; and each $N_{12}(c)$ credit is
attached to its own blocker incidence. If one flip improves
several such incidences, \Cref{lem:cv-fresh} counts the corresponding sum
of distinct metric-discount increments.
\end{proof}

\subsubsection{Small multiplicities: exact finite verification}

\begin{lemma}[Finite low-multiplicity estimate]\label{lem:cv-certificate}
For every regular colour \(c\) with
\(\mathsf A_c=1\) and \(k_c\in\set{1,2}\), every component state encoded
above and every $(\gamma_{\rm gain},\gamma_{\rm loss})$ in the rectangle
\eqref{eq:cv-gamma-box}
satisfy
\begin{equation}\label{eq:cv-certificate}
\mathcal H_c+\gamma_{\rm loss}N_{11}(c)
       -\gamma_{\rm gain}N_{12}(c)
 \le-1+k_c\Lambda_{\rm low}.
\end{equation}
The maximum excess of the left-hand side over the right-hand side of
\eqref{eq:cv-certificate}, over the stated records and coefficient pairs,
is $0$, attained at the balancing pair.
\end{lemma}

\begin{proof}[Computer-assisted proof]
By \Cref{lem:cv-balancing-domination}, it is enough to check the single
coefficient pair
\[
 (\gamma_{\rm gain},\gamma_{\rm loss})
   =\left(\frac{19}{500},\frac{33}{250}\right),
 \qquad \Lambda_{\rm low}=\frac{226}{125}.
\]
Exact rational enumeration gives maximum excess zero over all $400$
structural states and $400$ permitted maximizer choices for $k_c=1$, and
all $193600$ structural states and $196964$ permitted maximizer choices for
$k_c=2$.  These include every admissible availability and feasibility
pattern, repeated-component marker, truncated size, and choice among tied
maximum components.
For example, equality at the balancing pair is attained by the one-index
record $a_1=2$, $b_1=1$ with all four availability and feasibility flags
equal to $1$: then $\mathcal H_c=423/500$, $N_{11}(c)=0$, and
$N_{12}(c)=1$.

Every component configuration maps to one of these records: order its listed
neighbour incidences as in \Cref{def:cv-component-records}, write $0$ after the
first listed incidence of a repeated component, truncate every size at $7$, and retain
its actual availability and feasibility flags.  The permitted records use
one state at size $0$, two at size $1$, and three at each size
$2,\ldots,7$, exactly as justified in that definition.  The enumeration
also allows some records that no graph realizes; this only enlarges the
maximum and is therefore conservative.

The companion Lean development proves this finite estimate and the
encoding of the actual component configurations in
\nolinkurl{ZeroFreeness/Coupling/CV/Certificate.lean} and its dependent modules;
the finite arithmetic is checked by Lean's kernel without importing
Python output~\cite{Shi2026Formalization}.
As an independent check, the Python verifier
\nolinkurl{anc/python/soft_cv_certificate_verifier.py} checks the larger
eleven-pair test specified in \texttt{anc/CERTIFICATE\_SPEC.md}.
This includes the balancing pair above, every permitted maximum-component
tie, and an equality witness.  The Python source and
frozen certificate are covered by \texttt{anc/SHA256SUMS}.
For reproducibility, the exhaustive checks are run from the repository root with
\begin{verbatim}
python3 anc/python/soft_cv_certificate_verifier.py --check
\end{verbatim}
and the reported state counts are part of the certificate specification.
\end{proof}

\begin{lemma}[High multiplicity and missing colours]\label{lem:cv-high}
Set
\[
 \Lambda_{\rm bulk}:=\tfrac43+P_2+\tfrac43P_4=\tfrac{1331}{750}.
\]
For a root-allowed regular colour \(c\) with $k_c\ge3$ active neighbours,
\[
 \mathcal H_c+\gamma_{\rm loss}N_{11}(c)
 \le-1+\Lambda_{\rm bulk}k_c.
\]
For a regular colour \(c\) absent from the active root list,
\[
 \mathcal H_c+\gamma_{\rm loss}k_c
 \le \frac{39743}{27000}k_c\le\Lambda_{\rm bulk}k_c.
\]
The second inequality is strict when $k_c>0$.
\end{lemma}

\begin{proof}
The root-containing components have size at least $k_c+1\ge4$, so in
\eqref{eq:cv-charge} each root/off-root term is at most
\(\max_{r\ge4}(r-2)P_r=2P_4\), giving \(4P_4\) for the two sides. For the
contribution at one regular neighbour,
\(\max_{1\le r,s\le7}
  \bigl(rP_r+sP_s-\min\set{P_r,P_s}\bigr)=1+P_2\).
Indeed, $\max_j jP_j=1$ and $\max_j(j-1)P_j=P_2$; subtracting the smaller
of \(P_r,P_s\) leaves one term of each type.  Equality occurs at
\((r,s)=(2,1)\).  A location receives the geometric charge
$\gamma_{\rm loss}\le P_2$ only when both records are singletons, in
which case its component term is at most $1$. Either way the location
contributes at most $1+P_2$. Hence the total is at most
$4P_4+k_c(1+P_2)$, and
\[
 4P_4+k_c(1+P_2)\le-1+\Lambda_{\rm bulk}k_c,
\]
with slack $(k_c-3)(1+4P_4)/3\ge0$; the
same bound in the aggregate-parameter hard setting is
\cite[Lemma~4.6]{CarlsonVigoda2024}. If the colour is missing from the
active root list there are no root-containing moves, and the exact
per-neighbour maximum is $1+P_2$ by the same calculation, matching
\cite[Lemma~4.7]{CarlsonVigoda2024}; adding the largest geometric
correction $\gamma_{\rm loss}\le799/5400$ gives
$39743/27000<\Lambda_{\rm bulk}$.
\end{proof}

\begin{lemma}[The two root colours]\label{lem:cv-special}
Fix $c_0\in\set{a,b}$.  For the two root colours, define
\[
\begin{aligned}
 k_a&:=\bigl|\{u\in N_{G^\tau}(v):X_u=Y_u=a,\ vu\in E_Y\}\bigr|,\\
 k_b&:=\bigl|\{u\in N_{G^\tau}(v):X_u=Y_u=b,\ vu\in E_X\}\bigr|.
\end{aligned}
\]
Thus each multiplicity is counted on the side where the corresponding
root-containing move starts.  The cost \(C_{c_0}(F_X,F_Y)\) defined above for the corresponding
$\set{a,b}$-component moves satisfies
\[
 C_{c_0}(F_X,F_Y)\le-\one{c_0\text{ is root-allowed}}
              +\Lambda_{\rm bulk}k_{c_0}.
\]
\end{lemma}

\begin{proof}
If $c_0$ is unavailable there is no root-containing move and the off-root
components meet disjoint sets of the \(k_{c_0}\) neighbours.  After the
common denominator \(nq\) is cleared, each contributes at most one, so
\(C_{c_0}(F_X,F_Y)\le k_{c_0}\le\Lambda_{\rm bulk}k_{c_0}\). Let \(c_0\) be
available. If \(k_{c_0}=0\), the singleton root move paired with holding
probability contributes \(-1\). If \(k_{c_0}\ge2\), an infeasible root move leaves
cost at most \(k_{c_0}\le\Lambda_{\rm bulk}k_{c_0}-1\); a feasible one paired with
holding probability has total cost at most
\[
  k_{c_0}+\max_j(j-2)P_j=k_{c_0}+\tfrac{22}{125},
\]
including the off-root component moves, and
$2\Lambda_{\rm bulk}-1-(2+\tfrac{22}{125})
=\tfrac{28}{75}>0$ settles the weakest case \(k_{c_0}=2\). If \(k_{c_0}=1\), let the unique
component left after deleting $v$ from the root-containing component have
size $s$.  By the block classification in
\Cref{lem:cv-move-partition}, on the opposite side it either avoids \(v\)
and is the unique off-root partner in this family, or it also reaches \(v\)
and belongs to the other root-colour move, leaving no off-root copy here.
In the first case, because the target colour is root-allowed, the
root-containing flip is feasible exactly when its sole off-root component
is feasible.  Thus an infeasible root move contributes $0$. A feasible root move with no
feasible off-root partner is paired with holding probability at cost
$\max_s(s-1)P_{s+1}=\tfrac{22}{125}$; and a feasible pair costs
$sP_s-(s+1)P_{s+1}\le\tfrac{44}{125}$. All are below
$\Lambda_{\rm bulk}-1$, with slack at least $\tfrac{317}{750}$.
\end{proof}

\subsection{Averaging and degree bounds}

\subsubsection{Averaging over active multiplicities}
A physical regular root colour \(c\) of multiplicity \(m_c\ge3\) has random
active multiplicity \(K_c\sim\mathrm{Bin}(m_c,\theta)\).  Thus \(K_c\)
denotes the random multiplicity before the activations are exposed, whereas
\(k_c=\abs{\Gamma_c}\) denotes its value for a fixed active outcome.  In the
root-allowed case, the corrected per-colour contribution for
\(1\le K_c\le2\) is
bounded by \eqref{eq:cv-certificate}, whose positive coefficient is at most
\[
 \overline\Lambda_{\rm low}:=
 2-P_2+\tfrac{799}{5400}=\tfrac{49247}{27000};
\]
for \(K_c\ge3\) use $\Lambda_{\rm bulk}$ from \Cref{lem:cv-high} and discard
the negative credit. In the root-missing case the rate is
\(39743/27000<\Lambda_{\rm bulk}\) for every \(K_c\). After the allowed case's
$-1$ baseline is separated,
the expected positive charge per
\emph{physical} neighbour is therefore at most
\[
 \Lambda_{\rm bulk}\theta+
 (\overline\Lambda_{\rm low}-\Lambda_{\rm bulk})
 \theta x^{m_c-2}\bigl(1+(m_c-2)\theta\bigr)
 \le\Lambda_{\rm bulk},
\]
since \(\theta x^{m_c-3}(1+(m_c-2)\theta)\le2\) and
$3\Lambda_{\rm bulk}-2\overline\Lambda_{\rm low}
=22627/13500>0$.  Thus both availability cases cost at most
$\Lambda_{\rm bulk}$ per physical neighbour after the baseline is
separated. The baseline is accounted for in
\Cref{lem:cv-assembly}.

\subsubsection{Degree budget}
Let
\[
 L:=\sum_{\substack{c\notin T\\m_c\in\set{1,2}}}m_c,\qquad
 H:=\deg_{G^\tau}(v)-L,\qquad
 B_v^{\partial}:=\sum_{c\in\colours}\bcount_v^\tau(c).
\]
Then $L+H+B_v^{\partial}\le\Deg$ and, because $L_x\le L$,
\begin{equation}\label{eq:cv-resource}
 (L-L_x)+H+B_v^{\partial}\le\Deg-L_x.
\end{equation}
Every distinct colour missing from the active root list is witnessed by an
active pinned-neighbour constraint forbidding that colour.  Choosing one witness
for each missing colour is injective, so their number is at most the number
of active root pinned-neighbour constraints, and hence at most \(B_v^{\partial}\).

\begin{lemma}[The drift envelope]\label{lem:cv-assembly}
Let $\mathcal D$ be the expected one-step drift of $d_x^{G,\tau}$ under the coupling
above. Then
\begin{equation}\label{eq:cv-envelope}
  nq\,\mathcal D\;\le\;-q
      +\theta\Lambda_{\rm low}L_x
      +\Lambda_{\rm bulk}(\Deg-L_x).
\end{equation}
\end{lemma}

\begin{proof}
Apply the drift decomposition \eqref{eq:cv-drift-by-colour}.  Its regular-colour
summand is exactly the expression controlled by
\Cref{lem:cv-certificate,lem:cv-high}; the root-colour summand has no
geometric correction because the input metric discounts only regular
neighbours.

Condition on the root pinned-neighbour constraint coins; they determine the active
root list and are independent of all root-edge and blocker bits.
Every colour \(c\) contributes a baseline: an available colour contributes $-1$
--- through \Cref{lem:cv-completion-cost} when
it has no active neighbour, through \eqref{eq:cv-certificate}
when $1\le k_c\le2$, through \Cref{lem:cv-high} when $k_c\ge3$, and
through \Cref{lem:cv-special} for the two target colours --- while a
missing colour contributes $0$.  The witness injection above charges each
missing baseline to a distinct root pinned-neighbour constraint.

For each allowed low-multiplicity colour \(c\), root availability is independent of
the root-edge bits, so
$\E[k_c\one{\text{allowed}}]=\theta m_cx^{\bcount_v^\tau(c)}$, and each active
neighbour carries at most $\Lambda_{\rm low}$ by
\eqref{eq:cv-certificate}.  Their total is at most
$\theta\Lambda_{\rm low}L_x$. Neighbours of \emph{unavailable} low colours have expected
count
\[
 \theta\sum_{\substack{c\notin T\\m_c\in\set{1,2}}}
      m_c(1-x^{\bcount_v^\tau(c)})\le L-L_x
\]
and rate at
most $39743/27000<\Lambda_{\rm bulk}$ each.  Physical
high-multiplicity regular neighbours are charged at most
$\Lambda_{\rm bulk}$ each by the binomial averaging above, in both
availability cases; root-colour neighbours at most $\Lambda_{\rm bulk}$
each by \Cref{lem:cv-special}.  Finally, missing baselines cost at most one
per pinned-neighbour constraint, and $1<\Lambda_{\rm bulk}$. Hence
\[
  nq\,\mathcal D\le-q+\theta\Lambda_{\rm low}L_x
    +\Lambda_{\rm bulk}\bigl[(L-L_x)+H+B_v^{\partial}\bigr],
\]
and \eqref{eq:cv-resource} proves \eqref{eq:cv-envelope}.
\end{proof}

\subsection{Numerical bound and contraction}
\begin{lemma}[Final numerical inequality]\label{lem:cv-scalar}
For $\varrho\ge1809/1000$, $\Deg\ge125$, and
$\theta,\phi\in[0,1]$, interpret $\phi_x=\phi$ in
\eqref{eq:cv-gammas}. Then
\[
  -\varrho+\Lambda_{\rm bulk}(1-\phi)
    +\theta\phi\,\max\set{\Lambda_1,\Lambda_2,\Lambda_3}
  \;\le\;-\frac{59}{226125},
\]
where the rates in \eqref{eq:cv-lambda-rates} are evaluated at the actual
coefficients \eqref{eq:cv-gammas}; in particular
\(\Lambda_3=226/125\).
\end{lemma}

\begin{proof}
Write $\eta_0=17/200$, $y=\theta\phi$, and
$C_\Deg=1+P_2+2/\Deg$.  For \(j=1,2,3\), set
\[
 F_j:=-\varrho+\Lambda_{\rm bulk}(1-\phi)+y\Lambda_j.
\]
For the first case this expression is
\[
 F_1=-\varrho+\Lambda_{\rm bulk}(1-\phi)
       +y(2-P_2+\eta_0)
       +\frac{\eta_0}{\varrho}y(C_\Deg-P_2y).
\]
For fixed $y$ we have $y\le\phi$, so $\phi=y$ maximizes $F_1$. The resulting
function decreases in $\varrho$ and as $\Deg$ grows. At
$(\varrho,\Deg)=(1809/1000,125)$ its derivative in $y$ is decreasing in
$y$ and at $y=1$ equals $\tfrac{34097}{1809000}>0$. Thus the first case
is maximized at that corner and $y=\phi=\theta=1$.

For the second case,
\[
 \Lambda_2=2-P_3-\eta_0+
    \frac{\eta_0}{\varrho}\left(1+\frac2\Deg\right)>0.
\]
First maximize in $\theta$, and then in the affine variable $\phi$, to obtain
\[
 F_2\le-\varrho+\max\set{\Lambda_{\rm bulk},\Lambda_2}.
\]
Both $-\varrho+\Lambda_{\rm bulk}$ and $-\varrho+\Lambda_2$ decrease in $\varrho$,
and the latter also decreases as $\Deg$ grows. Hence their maximum occurs
at $(\varrho,\Deg)=(1809/1000,125)$; there
$\Lambda_2-\Lambda_{\rm bulk}=\tfrac{61637}{1809000}>0$, so
$\phi=\theta=1$ is the
maximizer at the worst corner. Finally
$\Lambda_3=226/125>\Lambda_{\rm bulk}$, and the same affine maximization
gives $F_3\le-\varrho+\Lambda_3$.

Exact evaluation gives the three gaps
\[
 \frac{59}{226125},\qquad\frac{59}{226125},\qquad\frac1{1000},
\]
whose minimum is \(\delta_{\rm CV}\).  The \texttt{scalar\_closure} certificate in
\nolinkurl{anc/python/soft_cv_certificate_verifier.py} verifies all three
corner values.
\end{proof}

\begin{theorem}[Geometric soft-CV contraction]\label{thm:cv-contraction}
Let \(q,\Deg\) be positive integers with \(\Deg\ge125\) and
\(q\ge1.809\,\Deg\). For every \(G\in\Gdeg\) and pinning \(\tau\) with
\(n=\lvert V^\tau\rvert\ge1\), every \(x\in(0,1]\), and every adjacent pair
\(X,Y\in\colours^{V^\tau}\),
there is a coupling \((X_1,Y_1)\) of
\(K_{x,\mathrm{CV}}^{G,\tau}(X)\) and \(K_{x,\mathrm{CV}}^{G,\tau}(Y)\) such that
\begin{equation}\label{eq:cv-contraction}
  nq\Bigl(\E[d_x^{G,\tau}(X_1,Y_1)]-d_x^{G,\tau}(X,Y)\Bigr)
  \le-\delta_{\rm CV}\Deg.
\end{equation}
\end{theorem}

\begin{proof}
Use the common-coin coupling of
\Cref{lem:coupled-activation-root-local} and the conditional component
coupling of \Cref{def:cv-component-matching,lem:cv-completion}.
The drift envelope \eqref{eq:cv-envelope}, which combines the component
charges and geometric-discount estimates, applies to this coupling.
\Cref{lem:cv-scalar} bounds its right-hand side by
$-\delta_{\rm CV}\Deg$.  At $x=1$ the metric is Hamming and the active
graph is empty.  The endpoint $x=0$ is handled below by continuity of the
normalized child laws, not by assigning values to $0^0$ in the active
representation.
\end{proof}

\subsection{From contraction to coupling independence}

It remains to compare each child kernel with the common middle kernel and
apply stationary comparison in the child metric.

\begin{lemma}[Soft-CV child--middle row discrepancy]
\label{lem:cv-child-middle}
Let \(G\in\Gdeg\), let \(\tau\) be a pinning, let \(r\in V^\tau\),
let \(c\in\colours\), and fix \(x\in(0,1]\). Define the middle kernel by
\[
  K_{{\rm mid},x}:=K_{x,\mathrm{CV}}^{G-r,\tau}.
\]
Its free graph is \(G^\tau-r\), with the original boundary counts restricted
to \(V^\tau\setminus\{r\}\). Thus \(r\) and all its incident interactions
are deleted without adding the pinned-neighbour constraints created by
pinning \(r=c\). If the common vertex set
\(V^\tau\setminus\{r\}\) has \(n_0\ge1\) vertices, then for every configuration
\(\sigma\in\colours^{V^\tau\setminus\{r\}}\),
\[
  \WHam\bigl(K_{x,\mathrm{CV}}^{G,\tau^{r=c}}(\sigma),
             K_{{\rm mid},x}(\sigma)\bigr)
       \le\frac{(1-x)\Deg}{n_0q}
       \le\frac{\Deg}{n_0q}.
\]
If $q\ge1.809\,\Deg$, the same bound holds in the child metric
$d_x^{G,\tau^{r=c}}\le\Ham$.
\end{lemma}

\begin{proof}
The child differs from the middle instance by at most \(\Deg\) pinned-neighbour
constraints.  Apply the common one-boundary estimate
\Cref{lem:boundary-sensitivity} with \(\rho_j=P_j\).  The CV parameters
satisfy \(\sup_j jP_j=1\), so the Hamming bound is
\((1-x)\Deg/(n_0q)\le\Deg/(n_0q)\).  Under the additional parameter
assumption, \Cref{lem:cv-metric} gives
\(d_x^{G,\tau^{r=c}}\le\Ham\), so the same coupling gives the
metric bound.
\end{proof}

\begin{proof}[Proof of \Cref{thm:cv-ci}]
Fix \(G\in\Gdeg\), a pinning \(\tau\), a root \(r\in V^\tau\), and
colours \(a,b\). If
\(a=b\), the laws coincide.  Let \(n_0=\abs{V^\tau\setminus\{r\}}\).  If
\(n_0=0\), both children are the empty point mass, so assume \(n_0\ge1\)
and first take \(x>0\).

Let \(K_{\rm mid}:=K_{x,\mathrm{CV}}^{G-r,\tau}\), whose stationary law is
\(\mu_{\rm mid}:=\mupin{G-r}{\tau}{x}\). For \(c\in\{a,b\}\), put
\[
  \mu_c:=\mupin{G}{\tau^{r=c}}{x}.
\]
Fix such a \(c\), and write \(d_c:=d_x^{G,\tau^{r=c}}\) for the child metric.
Then
\Cref{thm:cv-contraction,lem:cv-metric} gives, for adjacent \(X,Y\),
\[
\begin{aligned}
  W_{1,d_c}\!\left(
      K_{x,\mathrm{CV}}^{G,\tau^{r=c}}(X),
      K_{x,\mathrm{CV}}^{G,\tau^{r=c}}(Y)\right)
  &\le d_c(X,Y)-\frac{\delta_{\rm CV}\Deg}{n_0q}\\
  &\le\left(1-\frac{\delta_{\rm CV}\Deg}{n_0q}\right)
       d_c(X,Y).
\end{aligned}
\]
By \Cref{lem:cv-child-middle}, the child--middle row discrepancy is at most
\((1-x)\Deg/(n_0q)\) in this metric.  Stationary comparison with
contraction gap \(\delta_{\rm CV}\Deg/(n_0q)\) therefore yields
\[
  W_{1,d_c}(\mu_c,\mu_{\rm mid})
  \le\frac{1-x}{\delta_{\rm CV}}.
\]
Since \(d_c\ge m_0\Ham\),
\[
  \WHam(\mu_c,\mu_{\rm mid})
  \le\frac{1-x}{m_0\delta_{\rm CV}}.
\]
Applying this for \(c=a,b\) and using the triangle inequality through
\(\mu_{\rm mid}\) gives
\[
  \WHam(\mu_a,\mu_b)
  \le\frac{2(1-x)}{m_0\delta_{\rm CV}}
  \le\frac{2}{m_0\delta_{\rm CV}}.
\]

At \(x=0\), hard feasibility (\Cref{lem:hard-feas}, \(q\ge\Deg+1\))
makes both hard child partition sums positive, and
\Cref{lem:finite-endpoint-closure} passes the uniform bound to their
normalized hard limits. This proves \eqref{eq:cv-ci} also at \(x=0\).
\end{proof}
